\lecture{2}
\title{Learning-Based Approaches to Robotics}

\begin{document}

\subsection{Has AI Solved Robotics?}

As spoken about on
\href{https://x.com/chooi_jeq/status/2096064315115839904}{X},
GPT-6 Astra scored 95\% on a robot control task
in which robots are tasked to, e.g.,
``Pick up the red block from the table
and place it inside this bowl.''

In essence, we have these \textbf{Omni models}
that are currently capable of taking
audio, text, image, and video inputs
and producing audio, text, image, and video outputs.
The current trend is to add robotics to these lists.

\textbf{Question:}
If AI is solving robotics, then why get into robotics at all?
\begin{itemize}
    \item There is limitless potential for impact,
    and the state of the art is moving fast in \emph{every} field.
    \item Hardware can't scale (quite) as fast as software.
    \begin{itemize}
        \item (But exponentials can still happen\dots
        consider robots building robots, for example.)
    \end{itemize}
    \item In general, the gap between ``vibe coder'' and ``expert''
    is still nontrivially large.
\end{itemize}

\subsection{Reinforcement Learning}

One way to train robots is \textbf{reinforcement learning (RL)},
which consists of the following steps.
\begin{enumerate}
    \item Make a simulator.
    \item Define and write a reward function.
    \item Take gradients for optimization.
    \item Deploy in the real world.
\end{enumerate}
One downside of RL is that it requires a simulator, because:
\begin{itemize}
    \item It's much safer and more convenient to implement.
    \item It's much easier to define and compute a reward function
    (because the full state is perfectly known).
\end{itemize}
Another downside of RL is that it requires fine-tuning
the reward function, because \emph{``misalignment''}.

\subsection{Behavior Cloning}

Another way to train robots is \textbf{behavior cloning}---the point being
that humans do things pretty well, so we should just
demonstrate to robots what humans would do.

This can be done via \emph{teleoperation}.
Training becomes a supervised-learning objective;
no simulation or reward function required.

Naturally, one might also consider \emph{imitation-guided RL},
which may possibly yield ``superhuman'' performance.

\subsection{Training in Practice}

For \emph{single-task learning},
one usually uses roughly 200 human demonstrations
to produce a map from image readings into actions (visuomotor policy).
For \emph{multitask behavior cloning},
one usually uses a lot of robot and internet data
to produce a \emph{language-conditioned} visuomotor policy.

\begin{remark}
    And then some more things were said.
\end{remark}

\end{document}